\section{Consequences for Ambiguity}

Theorem~\ref{thm:trace-run} transports standard ambiguity classes to
backtracking traces.

\begin{corollary}[Unambiguous expressions]
If the position automaton for \(r\) is unambiguous, then for every word \(w\)
the trace tree \(\trace(r,w)\) has at most one successful leaf.
\end{corollary}

\begin{corollary}[Finite ambiguity]
If the position automaton for \(r\) is \(k\)-ambiguous, then every trace tree
\(\trace(r,w)\) has at most \(k\) successful leaves.
\end{corollary}

\begin{corollary}[Polynomial ambiguity growth]
If the position automaton for \(r\) has ambiguity growth bounded by a polynomial
\(p\), then
\[
  |\mathsf{SuccLeaves}(\trace(r,w))| \le p(|w|)
\]
for all \(w\).
\end{corollary}

These corollaries are deliberately phrased over successful leaves.  They do not
claim that the full tree is small.  A failing input can have many failing
branches even when successful ambiguity is low.  Nevertheless, successful
ambiguity is a semantic invariant that can be checked independently of a
particular engine implementation, and it provides a foundation for richer
analysis of failing prefixes.
